% alg-p3-03: 十字相乘法的"十字"图示
% 展示 ax^2 + bx + c 型因式分解的搜索结构
\documentclass[tikz,border=4pt]{standalone}
\usepackage{ctex}
\usepackage{amsmath}
\usetikzlibrary{arrows.meta, calc}
\begin{document}
\begin{tikzpicture}[scale=1.0, thick]

  % ---- 上方：x^2+bx+c 型（首项系数=1）----
  \begin{scope}[shift={(0,0)}]
    \node[font=\small, align=center] at (0, 1.5)
      {$x^2 + 5x + 6$（首项系数为 $1$）};

    % 十字线
    \draw[gray, thin] (-1.5,0) -- (1.5,0);    % 横线
    \draw[gray, thin] (0,-1.0) -- (0,1.0);    % 竖线

    % 四个角的数字
    \node[left, font=\normalsize]  at (-0.15, 0.5)  {$x$};
    \node[right, font=\normalsize] at ( 0.15, 0.5)  {$+2$};
    \node[left, font=\normalsize]  at (-0.15,-0.5)  {$x$};
    \node[right, font=\normalsize] at ( 0.15,-0.5)  {$+3$};

    % 对角线（交叉箭头，红色）
    \draw[-{Stealth[length=4pt]}, red, thick]
      (-0.8, 0.5) -- (0.8,-0.5);
    \draw[-{Stealth[length=4pt]}, red, thick]
      (-0.8,-0.5) -- (0.8, 0.5);

    % 交叉积说明
    \node[right, font=\footnotesize, red] at (1.6, 0) {$x{\cdot}3 + x{\cdot}2 = 5x$};

    % 左列相乘 → 首项，右列相乘 → 常数项
    \draw[|<->|, blue, thin] (-1.8, -0.6) -- (-1.8, 0.6);
    \node[left, font=\footnotesize, blue] at (-1.85, 0) {$x{\cdot}x=x^2$};

    \draw[|<->|, blue, thin] (1.0+0.7, -0.6) -- (1.0+0.7, 0.6);
    \node[right, font=\footnotesize, blue] at (1.0+0.75, 0) {$2{\cdot}3=6$};

    % 结论
    \node[font=\small] at (0, -1.5)
      {$\Rightarrow\; x^2+5x+6 = (x+2)(x+3)$};
  \end{scope}

  % ---- 下方：ax^2+bx+c 型（首项系数≠1）----
  \begin{scope}[shift={(0,-3.8)}]
    \node[font=\small, align=center] at (0, 1.5)
      {$2x^2 + 7x + 3$（首项系数 $\neq 1$）};

    % 十字线
    \draw[gray, thin] (-1.5,0) -- (1.5,0);
    \draw[gray, thin] (0,-1.0) -- (0,1.0);

    % 四个角
    \node[left, font=\normalsize]  at (-0.15, 0.5)  {$x$};
    \node[right, font=\normalsize] at ( 0.15, 0.5)  {$+3$};
    \node[left, font=\normalsize]  at (-0.15,-0.5)  {$2x$};
    \node[right, font=\normalsize] at ( 0.15,-0.5)  {$+1$};

    % 对角线
    \draw[-{Stealth[length=4pt]}, red, thick]
      (-0.8, 0.5) -- (0.8,-0.5);
    \draw[-{Stealth[length=4pt]}, red, thick]
      (-0.8,-0.5) -- (0.8, 0.5);

    % 交叉积
    \node[right, font=\footnotesize, red] at (1.6, 0)
      {$x{\cdot}1 + 2x{\cdot}3 = 7x$\;✓};

    % 左列右列
    \draw[|<->|, blue, thin] (-1.8, -0.6) -- (-1.8, 0.6);
    \node[left, font=\footnotesize, blue] at (-1.85, 0) {$x{\cdot}2x=2x^2$};

    \draw[|<->|, blue, thin] (1.0+0.7, -0.6) -- (1.0+0.7, 0.6);
    \node[right, font=\footnotesize, blue] at (1.0+0.75, 0) {$3{\cdot}1=3$};

    % 结论
    \node[font=\small] at (0,-1.5)
      {$\Rightarrow\; 2x^2+7x+3 = (x+3)(2x+1)$};
  \end{scope}

\end{tikzpicture}
\end{document}
